\documentclass[preview]{standalone}

\usepackage{amsmath}
\usepackage{amssymb}
\usepackage{stellar}
\usepackage{definitions}

\begin{document}

\id{psicologia-ereditarieta-comportamento}
\genpage

\section{Ereditarietà e comportamento}

\begin{snippet}{ereditarieta-comportamento-introduzione}
    Nello studio psicologico del comportamento, una questione centrale riguarda
    il rapporto tra \emph{natura} e \emph{cultura}, cioè tra
    \emph{ereditarietà} e \emph{ambiente}. L'ereditarietà rimanda al corredo
    biologico ricevuto dai genitori, mentre l'ambiente comprende le condizioni
    fisiche, sociali e culturali in cui l'individuo cresce e agisce.
\end{snippet}

\begin{snippetexample}{comportamento-aggressivo-example}{Comportamento aggressivo}
    Consideriamo il comportamento aggressivo. Una spiegazione ereditaria
    potrebbe sostenere che alcuni individui siano aggressivi perché possiedono
    una predisposizione biologica trasmessa dai genitori. Una spiegazione
    ambientale, invece, potrebbe sostenere che il grado di aggressività aumenti
    in funzione delle esperienze vissute, dei modelli osservati e delle
    condizioni sociali in cui l'individuo cresce.
\end{snippetexample}

\begin{snippetdefinition}{selezione-naturale-definition}{Selezione naturale}
    La \emph{selezione naturale} è il processo evolutivo per cui, in un dato
    ambiente, gli individui dotati di caratteristiche più favorevoli alla
    sopravvivenza e alla riproduzione tendono a lasciare più discendenti
    rispetto agli altri membri della stessa specie.
\end{snippetdefinition}

\begin{snippet}{darwin-selezione-naturale}
    Charles Darwin presenta la teoria dell'evoluzione per selezione naturale
    ne \emph{L'origine delle specie}, pubblicato nel 1859. La varietà delle
    specie viene spiegata attraverso l'accumulo, nel corso delle generazioni,
    di tratti che offrono un vantaggio riproduttivo in un determinato ambiente.
\end{snippet}

\includesnpt[width=58\%,src=/snippet/static-psychology/natural-selection-process.jpg]{centered-img}

\begin{snippetdefinition}{genotipo-definition}{Genotipo}
    Il \emph{genotipo} è la struttura genetica ereditata da un individuo al
    momento del concepimento.
\end{snippetdefinition}

\begin{snippetdefinition}{fenotipo-definition}{Fenotipo}
    Il \emph{fenotipo} è l'insieme delle caratteristiche osservabili di un
    individuo, compresi il suo aspetto fisico e il repertorio dei suoi
    comportamenti.
\end{snippetdefinition}

\begin{snippetdefinition}{vantaggio-selettivo-definition}{Vantaggio selettivo}
    Un \emph{vantaggio selettivo} è una caratteristica che, in un determinato
    ambiente, aumenta la probabilità che un individuo sopravviva e si riproduca.
\end{snippetdefinition}

\begin{snippetexample}{giraffa-vantaggio-selettivo-example}{Giraffa e vantaggio selettivo}
    Una giraffa eredita dai genitori un certo genotipo. In un particolare
    ambiente, tale genotipo contribuisce allo sviluppo di un fenotipo, per
    esempio un collo lungo e la capacità di raggiungere foglie poste in alto.
    Se le foglie basse non sono sufficienti per l'intera popolazione, una
    giraffa dal collo lungo possiede un vantaggio selettivo perché può accedere
    a una fonte di cibo meno disponibile agli altri individui.
\end{snippetexample}

\begin{snippet}{adattamenti-specie-umana}
    Nell'evoluzione della specie umana, la selezione naturale ha favorito due
    adattamenti particolarmente rilevanti: il \emph{bipedismo} e
    l'\emph{encefalizzazione}. Il bipedismo ha modificato il rapporto del corpo
    con l'ambiente, mentre l'aumento delle dimensioni e della complessità del
    cervello ha reso possibili capacità più articolate di pensiero, memoria,
    pianificazione e ragionamento. Un'ulteriore tappa fondamentale è stata la
    comparsa del \emph{linguaggio}, alla base dell'evoluzione culturale.
\end{snippet}

\begin{snippetdefinition}{genetica-definition}{Genetica}
    La \emph{genetica} è lo studio dei geni e dei meccanismi dell'ereditarietà,
    cioè della trasmissione biologica dei caratteri dai genitori ai figli.
\end{snippetdefinition}

\begin{snippetdefinition}{genetica-comportamento-umano-definition}{Genetica del comportamento umano}
    La \emph{genetica del comportamento umano} è l'area di studio che esamina
    il contributo dei fattori genetici alle differenze individuali nel
    comportamento, nelle capacità cognitive e nelle caratteristiche della
    personalità.
\end{snippetdefinition}

\begin{snippetdefinition}{gene-definition}{Gene}
    Un \emph{gene} è un segmento di DNA coinvolto nella trasmissione dei
    caratteri ereditari e nella produzione di proteine.
\end{snippetdefinition}

\begin{snippetdefinition}{carattere-poligenico-definition}{Carattere poligenico}
    Un \emph{carattere poligenico} è un tratto alla cui manifestazione
    contribuiscono più geni.
\end{snippetdefinition}

\begin{snippetdefinition}{ereditabilita-definition}{Ereditabilità}
    L'\emph{ereditabilità} è una stima dell'influenza relativa della componente
    genetica, rispetto a quella ambientale, sulle differenze individuali in un
    tratto o comportamento. Si misura su una scala da \(0\) a \(1\): valori
    vicini a \(1\) indicano una maggiore associazione tra differenze genetiche
    e differenze osservate nel tratto considerato.
\end{snippetdefinition}

\begin{snippetdefinition}{studi-adozioni-definition}{Studi sulle adozioni}
    Gli \emph{studi sulle adozioni} confrontano le caratteristiche di bambini
    cresciuti in famiglie adottive con quelle dei genitori biologici e dei
    genitori adottivi, per valutare il peso relativo di ereditarietà e ambiente.
\end{snippetdefinition}

\begin{snippetdefinition}{studi-gemelli-definition}{Studi sui gemelli}
    Gli \emph{studi sui gemelli} confrontano la somiglianza tra gemelli
    monozigoti e gemelli dizigoti rispetto a un tratto o comportamento. Poiché
    i gemelli monozigoti condividono una quota genetica maggiore, il confronto
    permette di stimare il contributo dell'ereditarietà.
\end{snippetdefinition}

\begin{snippet}{interazione-geni-ambiente}
    Geni e ambiente non agiscono come cause isolate. Il fenotipo emerge dalla
    loro interazione: il corredo genetico fornisce possibilità biologiche, ma
    le esperienze e il contesto contribuiscono a orientare il modo in cui tali
    possibilità si manifestano nel comportamento.
\end{snippet}

\end{document}
